\begin{frame}[t]{Research record: biophysics}
\vspace{0pt}
{\fontsize{8.1}{9.8}\selectfont
{\color{Muted}PhD, University of Houston, with the Rice Center for Theoretical Biological Physics}\\[2pt]
\textbf{Eliaz}, Nedelec, Morrison, Levine \& Cheung. Insights from graph theory on the morphologies of actomyosin networks with multilinkers. \emph{Phys.\ Rev.\ E} 102:062420 (2020). First author.\\
Liman, \dots, \textbf{Eliaz}, \dots, Cheung. The role of the Arp2/3 complex in shaping the dynamics and structures of branched actomyosin networks. \emph{PNAS} 117:10825 (2020).\\
Li, Liman, \textbf{Eliaz} \& Cheung. Forecasting avalanches in branched actomyosin networks with network science and machine learning. \emph{J.\ Phys.\ Chem.\ B} 125:11591 (2021).\\
Zegarra, \dots, \textbf{Eliaz}, \dots. Impact of hydrodynamic interactions on protein folding rates depends on temperature. \emph{Phys.\ Rev.\ E} 97:032402 (2018).\\
Zhang, Nde, \textbf{Eliaz}, Jennings, Cieplak \& Cheung. CaXML: Chemistry-informed machine learning explains mutual changes between protein conformations and calcium ions in calcium-binding proteins using structural and topological features. \emph{Protein Science} 34:e70023 (2025).\\[4pt]
\textbf{Thesis}\\
\textbf{Eliaz}. Physics of self-assembly in complex matter. PhD thesis, University of Houston (2020); self-published in book form as \emph{Physics of Complex (Biological) Matter} (2023).\par}
\sources{Full author lists and positions: CLAIM-FENCE.md. Meeting abstracts (Biophysical Society, APS March Meeting 2018--2021) are not listed as papers.}
\note{[backup] PhD work at the University of Houston and Rice: graph theory of actomyosin network shape (PRE 2020, linker valency), Arp2/3 branching and avalanches (PNAS 2020, JPCB 2021), protein folding and calcium binding. Used on slide 7 (the journey) and slide 44, the declared prior for Test 2. It motivates Test 2 and is no evidence for it. Team results with co-authors.}
\end{frame}
